\section*{Chapter \ref{ch:g12:mountain-building} --- Building Mountains}

\begin{solution}{exo:g12:mountain-building:1}
(a) \omterm{def:g12:mountain-building:thrust}{thrust}; (b) \omterm{def:g12:mountain-building:ophiolite}{ophiolite}; (c) \omterm{def:g12:mountain-building:fold}{syncline}; (d) \omterm{def:g12:mountain-building:suture}{suture}; (e) \omterm{def:g12:mountain-building:fold}{anticline};
(f) \omterm{def:g12:mountain-building:nappe}{nappe}.
\end{solution}

\begin{solution}{exo:g12:mountain-building:2}
(1) The limbs dip away from the \omterm{def:g12:mountain-building:fold}{fold}'s axis: an \omterm{def:g12:mountain-building:fold}{anticline}. (2) The oldest
layers. (3) No: erosion follows the hard and the soft layers, not the shape
of the \omterm{def:g12:mountain-building:fold}{fold}; the core of an \omterm{def:g12:mountain-building:fold}{anticline} can be worn into a valley.
\end{solution}

\begin{solution}{exo:g12:mountain-building:3}
(1) $e = 1 - \cos\ang{30} = 1 - 0.866 = 0.134$, about \qty{13}{\percent}.
(2) $L_0 = 6.0/\cos\ang{30} \approx \qty{6.9}{km}$.
\end{solution}

\begin{solution}{exo:g12:mountain-building:4}
(1) A \omterm{def:g11:ocean-ridges:normal-fault}{normal fault}; the crust was stretched. (2) Yes: it is a \omterm{def:g12:mountain-building:thrust}{reverse fault}
dipping less than \ang{45}.
\end{solution}

\begin{solution}{exo:g12:mountain-building:5}
(1) Marine sediments, pillow basalts, sheeted \omterm{def:g11:igneous-rocks:intrusion}{dykes}, \omterm{def:g11:ocean-ridges:gabbro}{gabbro}, \omterm{def:g11:ocean-ridges:peridotite}{peridotite}.
(2) Lava forms pillows only when it cools in water. (3) An \omterm{def:g12:mountain-building:ophiolite}{ophiolite} is a
piece of the ocean that lay between two continents; when they collided,
slices of its floor were pushed onto their edges instead of sinking, so they
are found where the ocean closed, along the \omterm{def:g12:mountain-building:suture}{suture}.
\end{solution}

\begin{solution}{exo:g12:mountain-building:6}
(1) $L_0 = 2.1 + 1.8 + 2.4 + 1.6 = \qty{7.9}{km}$;
$e = (7.9 - 6.0)/7.9 \approx 0.24$, \qty{24}{\percent}. (2) $\cos\theta =
6.0/7.9 \approx 0.759$, $\theta \approx \ang{41}$.
\end{solution}

\begin{solution}{exo:g12:mountain-building:7}
(1) $45\cos\ang{10} \approx \qty{44.3}{km}$. (2) $45\sin\ang{10} \approx
\qty{7.8}{km}$. (3) Such \omterm{def:g12:mountain-building:fold}{folds} shorten by $1 - \cos\ang{12} \approx 0.0219$:
the belt must have been $44.3/0.0219 \approx \qty{2000}{km}$ wide before
folding.
\end{solution}

\begin{solution}{exo:g12:mountain-building:8}
(1) Superposition: read without the fault, the upper rock would be taken
for the younger. (2) $500 - 75 = 425$ million years. (3) $\qty{40000}{m} /
\qty{800}{m} = 50$ times.
\end{solution}

\begin{solution}{exo:g12:mountain-building:9}
(1) $(30.23 - 19.38) \times 111.2 = 10.85 \times 111.2 \approx
\qty{1207}{km}$. (2) $1207/14 \approx \qty{86}{km}$ per million years,
\qty{86}{mm} or about \qty{8.6}{cm} a year. (3) \qty{9}{m} a century is
\qty{9}{cm} a year: the two agree within about \qty{5}{\percent}.
\end{solution}

\begin{solution}{exo:g12:mountain-building:10}
(1) $H = 35/(1 - 0.40) \approx \qty{58}{km}$. (2) $\Delta H \approx
\qty{23.3}{km}$; it stands higher by $23.3 \times 550/3380 \approx
\qty{3.8}{km}$, and its extra \omterm{def:g12:continental-crust:root}{root} is $23.3 - 3.8 \approx \qty{19.5}{km}$ deep.
\end{solution}

\begin{solution}{exo:g12:mountain-building:11}
(1) The drawing multiplies $\tan\theta$ by 3: $\tan\theta = \tan\ang{40}/3
\approx 0.280$, $\theta \approx \ang{15.6}$. (2) True: $1 - \cos\ang{15.6}
\approx 0.037$, \qty{3.7}{\percent}; read on the drawing: $1 - \cos\ang{40}
\approx 0.234$, \qty{23}{\percent}. (3) A stretched vertical scale steepens
the dips and lengthens the layers, so it makes the \omterm{def:g12:mountain-building:shortening}{shortening} look several
times too large.
\end{solution}

\begin{solution}{exo:g12:mountain-building:12}
(1) $L_0 = 10/(1 - 0.317) \approx \qty{14.6}{km}$. (2) $1 - \cos\ang{57.5}
\approx 1 - 0.537 = 0.463$, \qty{46}{\percent}. (3) In the rounded \omterm{def:g12:mountain-building:fold}{fold}
only a short part of the layer dips at \ang{57.5}; near the hinges it is
almost flat, and flat parts are not shortened. In the straight-limbed \omterm{def:g12:mountain-building:fold}{fold}
the whole layer dips at \ang{57.5}.
\end{solution}

\begin{solution}{exo:g12:mountain-building:13}
(1) $\sqrt{16.6^2 + 41.8^2} \approx \qty{45.0}{mm/yr}$. (2) $41.8 - 23.5 =
\qty{18.3}{mm}$ a year, that is about \qty{18}{km} of north--south \omterm{def:g12:mountain-building:shortening}{shortening}
per million years north of Lhasa. (3) Lhasa itself has moved north, and the
earthquakes of the belt spread far north of the Himalaya across Asia: the
crust there is still breaking.
\end{solution}

\begin{solution}{exo:g12:mountain-building:14}
(1) Its dip is close to zero, so the \omterm{def:g12:mountain-building:shortening}{shortening} is the displacement,
\qty{45}{km}. (2) $45/(1 - \cos\ang{12}) = 45/0.0219 \approx \qty{2060}{km}$.
(3) Large convergences are absorbed mainly by \omterm{def:g12:mountain-building:thrust}{thrusts} and \omterm{def:g12:mountain-building:nappe}{nappes} and by tight
\omterm{def:g12:mountain-building:fold}{folds}; gentle \omterm{def:g12:mountain-building:fold}{folds} absorb very little.
\end{solution}

\begin{solution}{exo:g12:mountain-building:15}
(1) $\qty{8848.86}{m} = \qty{8848860}{mm}$ in $\num{5e7}$ years: about
\qty{0.18}{mm} a year. (2) $40/0.18 \approx 220$: the rise is more than two
hundred times slower than India's motion. (3) The limestone was laid down
below sea level, and erosion has removed the rocks that lay above it: it has
been lifted by more than its present height.
\end{solution}

\begin{solution}{pb:g12:mountain-building:1}
\textbf{1.} $2\pi \times 6371/360 \approx \qty{111.2}{km}$.

\textbf{2.} Between 50 million years ago (\ang{0.68}~S) and 45 (\ang{3.30}~N):
$50 - 5 \times 0.68/3.98 \approx 49$ million years ago.

\textbf{3.} $(19.38 - 0.68) \times 111.2 = 18.70 \times 111.2 \approx
\qty{2079}{km}$.

\textbf{4.} $2079/16 \approx \qty{130}{km}$ per million years: \qty{130}{mm}
a year, about \qty{13}{cm}.

\textbf{5.} 66 to 60: $7.28 \times 111.2/6 \approx 135$; 60 to 55: $4.86
\times 111.2/5 \approx 108$; 55 to 50: $6.56 \times 111.2/5 \approx 146$; 50
to 45: $3.98 \times 111.2/5 \approx 89$; 45 to 40: $2.29 \times 111.2/5
\approx 51$ (\unit{mm} a year). Fastest from 55 to 50 million years ago.

\textbf{6.} $(20.00 - 5.59) \times 111.2/40 = 14.41 \times 111.2/40 \approx
\qty{40}{mm}$ a year.

\textbf{7.} $40/130 \approx 0.31$: less than a third.

\textbf{8.} From the step 55--50 (146) to the step 50--45 (89): around 50
million years ago.

\textbf{9.} \omterm{def:g12:continental-crust:continental-crust}{Continental crust} is too light to sink: when it reached the
trench, \omterm{def:g11:subduction-zones:subduction}{subduction} could no longer pull India along freely, and the other
continent held it back.

\textbf{10.} $(20.00 + 0.68) \times 111.2 = 20.68 \times 111.2 \approx
\qty{2300}{km}$.

\textbf{11.} $(29.65 - 21.15) \times 111.2 = 8.50 \times 111.2 \approx
\qty{945}{km}$.

\textbf{12.} $2300 - 945 \approx \qty{1354}{km}$.

\textbf{13.} Today $(29.65 - 20.00) \times 111.2 \approx \qty{1073}{km}$;
50 million years ago $(21.15 + 0.68) \times 111.2 = 21.83 \times 111.2
\approx \qty{2427}{km}$ (and $2427 - 1073 = 1354$).

\textbf{14.} $e = 1354/2427 \approx 0.56$, \qty{56}{\percent}.

\textbf{15.} \omterm{def:g12:mountain-building:fold}{Folds}, and \omterm{def:g12:mountain-building:thrust}{reverse faults}: \omterm{def:g12:mountain-building:thrust}{thrusts} carrying \omterm{def:g12:mountain-building:nappe}{nappes}.

\textbf{16.} $H = 35 \times 2427/1073 \approx \qty{79}{km}$.

\textbf{17.} $\Delta H = 79 - 35 = \qty{44}{km}$; it would stand higher by
$44 \times 550/3380 \approx \qty{7.2}{km}$.

\textbf{18.} $\Delta H = 4.6 \times 3380/550 \approx \qty{28}{km}$, a crust
$35 + 28 = \qty{63}{km}$ thick, close to the \qty{70}{km}
measured.

\textbf{19.} The crust does not keep its cross-section: parts of Asia are
squeezed out eastward, and erosion removes rock from the top of the
mountains.

\textbf{20.} $2079 + 2300 \approx \qty{4380}{km}$: since the end of the
Cretaceous, India has travelled about \textbf{\qty{4400}{km} north}.
\end{solution}
